\item El termómetro de espíritu en vidrio, inventado en Florencia, Italia, alrededor de 1654, consiste de un tubo de líquido (el espíritu) que contiene un cierto número de esferas de vidrio sumergidas con masas ligeramente distintas (figura P14.41). A temperaturas suficientemente bajas, todas las esferas flotan, pero conforme se eleva la temperatura, las esferas se hunden una tras otra. El dispositivo es un burdo pero interesante instrumento para medir la temperatura. Suponga que el tubo está lleno con alcohol etílico, cuya densidad es $0.789\,45 \text{ g/cm}^3$ a $20.0^\circ\text{C}$ y disminuye a $0.780\,97 \text{ g/cm}^3$ a $30.0^\circ\text{C}$. (a) Suponga que una de las esferas tiene un radio de $1.000 \text{ cm}$ y está en equilibrio a la mitad del tubo a $20.0^\circ\text{C}$, determine su masa. (b) Cuando la temperatura se incrementa a $30.0^\circ\text{C}$, ¿qué masa debe tener una segunda esfera del mismo radio para estar en equilibrio en el punto medio? (c) A $30.0^\circ\text{C}$, la primera esfera ha caído al fondo del tubo. ¿Qué fuerza hacia arriba ejerce el fondo del tubo sobre esta esfera?
